\section{Designentscheidungen}

Bei der Umsetzung des High-Fidelity-Prototyps wurden mehrere zentrale Designentscheidungen getroffen.
Diese werden im Folgenden ausführlich begründet und in Bezug auf die Nutzer- und Aufgabenanalyse aus Meilenstein~1
sowie auf die dort durchgeführte Konkurrenzanalyse gesetzt.

\subsection{Karte als primäre Interaktionsoberfläche}

Die grundlegendste Designentscheidung des Projekts war die Wahl der Karte als zentrales Interface-Element.
Statt einer Listen- oder Dashboard-Ansicht steht beim Start der App sofort eine interaktive Karte Wiens im Vordergrund.

Diese Entscheidung leitet sich direkt aus der Kontextanalyse in M1 ab:
Die App wird primär mobil und \textit{unterwegs} verwendet, in öffentlichen Verkehrsmitteln, auf dem Weg zu einem Ort, oder spontan in einer Pause.
In diesem Nutzungskontext hat eine Karte einen entscheidenden Vorteil gegenüber einer abstrakten Liste:
Sie visualisiert Datenpunkte räumlich und erlaubt Nutzer:innen, sofort einen Bezug zur eigenen Umgebung herzustellen.
Das entspricht dem HCI-Prinzip der \textit{direkten Manipulation}: Nutzer:innen interagieren nicht mit abstrahierten Daten,
sondern mit ihrer tatsächlichen Stadt.

Aus der Konkurrenzanalyse (M1) ging hervor, dass mobile Open-Data-Anwendungen für Wien im App-Markt kaum vertreten sind.
Webbasierte Angebote wie \textit{Winterthur in Zahlen} oder der \textit{OECD Data Explorer} bieten zwar Kartenansichten,
sind aber nicht auf mobile, kontextbezogene Nutzung ausgelegt.
Die konsequente Kartenzentrierung füllt diese Lücke.

Zudem hat diese Version unserer Low-fidelity Prototypen bei den Evaluierungen, sowie bei der Präsentation am meisten Anhang gefunden.

\subsection{Farbcodierung der Kartenmarkierungen nach Kategorie}

Auf der Karte sind Pins in fünf Farben eingeteilt: Orange (Kultur), Blau (Verkehr), Grün (Umwelt),
Lila (Infrastruktur) und Grau (Basiskarte).
Ohne diese Unterscheidung wäre die Karte bei dieser Datendichte schwer lesbar. Alle Pins würden gleich aussehen
und Nutzer:innen müssten jeden einzelnen antippen, um seine Kategorie zu erfahren.

Die Farbcodierung folgt dem Gestalt-Prinzip der \textit{Ähnlichkeit}:
Elemente gleicher Farbe werden als zusammengehörig wahrgenommen, ohne dass eine Legende aktiv gelesen werden muss.
Damit wird die in der Aufgabenanalyse (M1) als wichtig eingestufte Aufgabe \enquote{Einträge kategorisieren} auf die
Kartenebene übertragen. Die Kategorisierung geschieht passiv und automatisch durch die visuelle Wahrnehmung.

Für Gelegenheitsnutzer:innen (Persona \textit{Luca Kunze}) ist dieser Ansatz besonders wichtig:
Sie haben laut M1 kein vertieftes Interesse an Daten und wollen ohne Vorkenntnisse schnell an Informationen gelangen.
Die intuitive Farbkodierung senkt die Einstiegshürde erheblich.
Gleichzeitig befriedigt sie technikaffine Studierende (Persona \textit{Benjamin Eberhart}),
die auf einen Blick die räumliche Verteilung bestimmter Datensatzkategorien erfassen wollen.

\subsection{Such-Bottom-Sheet statt separatem Such-Screen}

Die Suchfunktion ist als ausfahrbares Bottom-Sheet unterhalb der Karte realisiert und nicht als eigener Screen mit
Navigation und Zurück-Button.
Im zugeklappten Zustand ist das Suchfeld dauerhaft sichtbar. Wenn man Tippt fährt das Sheet auf und zeigt
frühere Suchanfragen oder Vorschläge.

Der Vorteil dieses Ansatzes liegt in der \textit{kontinuierlichen Kontextwahrnehmung}:
Die Karte bleibt im Hintergrund sichtbar und die Kamera animiert sich beim Auswählen eines Suchergebnisses direkt
auf den gefundenen Ort. Nutzer:innen sehen sofort, \textit{wo} sich das Ergebnis befindet.
Ein separater Such-Screen hätte diesen räumlichen Kontext unterbrochen.

Dieses Muster folgt etablierten Designkonventionen von Google Maps und Apple Maps und reduziert den kognitiven Aufwand,
der durch Bildschirmwechsel entsteht (\textit{Hick's Law}: je weniger Entscheidungsschritte,
desto schneller die Aufgabenerledigung).
Es adressiert die in M1 beschriebene Anforderung an \textit{schnelle, effiziente Navigation} bei kurzen Nutzungssessions.

\subsection{Verbindung zwischen Karte und Downloads-Screen}

Eine wichtige Interaktionsentscheidung ist die bidirektionale Verlinkung zwischen Kartenansicht und Downloads-Screen:
Wird auf der Karte ein Pin angetippt, zeigt das Detail-Modal den zugehörigen Datensatz mit einem Weiterleitungs-Button.
Umgekehrt können Nutzer:innen aus der Pinliste eines installierten Datensatzes direkt auf die Karte navigieren,
wobei die Kamera automatisch auf den gewählten Ort zoomt.

Diese Entscheidung spiegelt die Erkenntnis aus der Aufgabenanalyse (M1) wider,
dass Nutzer:innen unterschiedliche \textit{Einstiegspunkte} in die App haben:
Manche Nutzer suchen auf der Karte und wollen dann mehr über einen Datensatz erfahren.
Andere beginnen im Downloads-Screen und wollen einen bestimmten Ort auf der Karte erkunden.
Beide Wege müssen ohne Reibungsverluste möglich sein.

\subsection{Swipe-to-reveal für sekundäre Aktionen im Downloads-Screen}

Die drei Aktionen Lesezeichen setzen, Teilen und Löschen sind nicht als dauerhaft sichtbare Buttons implementiert,
sondern werden durch eine horizontale Wischgeste nach links freigelegt.

Diese Entscheidung basiert auf zwei Überlegungen:
Erstens verhindert sie \textit{visuelle Überladung} der Listeneinträge.
Jeder Datensatz zeigt nur die Hauptinformationen: Name, Beschreibung, Größe und Downloadstatus.
Sekundäre Aktionen bleiben verborgen und stören die Übersichtlichkeit nicht.
Zweitens schützt das Muster vor \textit{versehentlichen Aktionen}:
Das Löschen eines heruntergeladenen Datensatzes erfordert eine bewusste Wischbewegung und dann einen zusätzlichen Tap damit wird beunabsichtiges Löschen vermieden.

Das Swipe-Muster ist ein auf Android und iOS weit verbreitetes Interaktionsmuster
(z.\,B. E-Mail-Clients, Aufgaben-Apps) und entspricht daher dem HCI-Prinzip der
\textit{Konsistenz mit Nutzererwartungen}.
Berufliche Nutzer:innen wie \textit{Ralph Ebersbach} aus M1 können so effizient zwischen Entdecken und Verwalten wechseln.

\subsection{Tab-Trennung in \enquote{Entdecken} und \enquote{Installiert}}

Der Downloads-Screen unterscheidet explizit zwischen verfügbaren und bereits installierten Datensätzen.
Diese Trennung adressiert zwei konzeptuell verschiedene Nutzungsszenarien, die in der Aufgabenanalyse (M1) identifiziert wurden:

\textit{Erkunden und Entdecken} (häufig bei Content-Creator-Persona \textit{Leah Reiniger} und Student \textit{Benjamin Eberhart}):
Hier ist eine vollständige, übersichtliche Liste aller verfügbaren Datensätze gewünscht -- ohne Rauschen durch bereits bekannte Einträge.
\textit{Verwalten und Wiederverwenden} (häufig bei Persona \textit{Ralph Ebersbach}):
Hier werden nur installierte Datensätze benötigt, die für wiederkehrende Aufgaben bereitstehen.

Ohne diese Trennung würde die Liste mit zunehmender Anzahl installierter Datensätze schwerer navigierbar.
Die Tab-Struktur hält beide Ansichten sauber und fokussiert.

\subsection{Offline-Verfügbarkeit als Kernfunktion}

Die Möglichkeit, Datensätze lokal herunterzuladen, ist keine optionale Erweiterung, sondern ein zentrales Feature der App.
Diese Entscheidung ergibt sich unmittelbar aus der Kontextanalyse in M1:
Die primäre Nutzungssituation ist mobil -- in der U-Bahn, auf dem Fahrrad, in Grünflächen -- und damit oft mit
instabiler oder fehlender Internetverbindung verbunden.

Die in M1 als potenzielles Problem identifizierten \enquote{langen Ladezeiten bei schwacher Verbindung}
werden durch Offline-Karten strukturell gelöst: Einmal heruntergeladen, funktionieren alle Pins und zugehörigen
Informationen vollständig ohne Netzanbindung.
Dies ist ein Beispiel für \textit{proaktives Fehlermanagement} im Sinne von Nielsens Heuristiken --
das Problem wird verhindert, bevor es entsteht, anstatt nach dem Auftreten behandelt zu werden.

\subsection{Barrierefreiheit und Darstellungsanpassung}

Unter \enquote{Einstellungen} bietet die App explizite Optionen für große Schrift, hohen Kontrast sowie
die Wahl zwischen hellem, dunklem und systemfolgendem Farbschema.

Die Entscheidung, Barrierefreiheit als dedizierte Einstellungskategorie zu gestalten und nicht nur auf
Betriebssystemebene zu delegieren, folgt aus der Diversität der identifizierten Nutzergruppen (M1):
Von technikaffinen Studierenden bis zu Gelegenheitsnutzer:innen unterschiedlicher Altersgruppen
nutzen die App sehr verschiedene Personen mit unterschiedlichen Anforderungen an Lesbarkeit und Kontrast.
Durch die In-App-Einstellung bleibt die Anpassung direkt erreichbar, ohne die Systemeinstellungen verlassen zu müssen.
